\pdfvariable trailerid {[<C2882026090200000000000000000000><C2882026090200000000000000000000>]}
\documentclass[10pt]{article}
\usepackage[margin=0.72in]{geometry}
\usepackage{amsmath,amssymb,amsthm,booktabs,microtype,hyperref}
\hypersetup{colorlinks=true,linkcolor=blue,urlcolor=blue}
\providecommand{\CRevisionRound}{2}
\newtheorem{theorem}{Theorem}
\newtheorem{lemma}{Lemma}
\newcommand{\R}{\mathbb R}
\newcommand{\D}{\mathrm d}
\newcommand{\erfc}{\operatorname{erfc}}
\title{Resolvent, Scattering, and Heat Dynamics\\
of a One-Dimensional Delta Interaction}
\author{Route-A source-local certificate HCS-C288}
\date{2 September 2026}
\begin{document}
\maketitle

\begin{abstract}
We give a convention-complete solution of the one-center delta interaction
on the real line in units $\hbar=2m=1$.
\ifnum\CRevisionRound=0
The closed quadratic form yields the exact continuity and derivative-jump
domain, and a scalar interface equation gives the rank-one resolvent.
\fi
\ifnum\CRevisionRound>0
Its pole produces exactly one attractive bound state; the remaining spectrum,
left--right scattering, and odd/even channels are closed for every real
coupling, including threshold and free faces.
\fi
\ifnum\CRevisionRound>1
Laplace inversion gives the heat kernel and an integrable relative diagonal
with explicit trace.  Independent exact and high-precision evidence audits
all constants without replacing the analytic proof.  The model has a natural
self-adjoint quantization but no arithmetic-orbit or target-determinant
bridge.
\fi
\end{abstract}

\section{The frozen singular Hamiltonian}
For $\alpha\in\R$, consider the form
\begin{equation}\label{eq:form}
 q_\alpha[\psi]=\int_\R |\psi'(x)|^2\,\D x
    +\alpha|\psi(0)|^2,
 \qquad \mathcal D(q_\alpha)=H^1(\R).
\end{equation}
The one-dimensional trace inequality makes point evaluation infinitesimally
form bounded relative to the Dirichlet integral.  Thus~\eqref{eq:form} is
closed and lower bounded and defines one self-adjoint operator $H_\alpha$.
Integration by parts identifies its domain:
\begin{equation}\label{eq:domain}
 \begin{split}
 \mathcal D(H_\alpha)=\{\psi\in H^2(\R\setminus\{0\}):{}&
 \psi(0+)=\psi(0-),\\[-2pt]
 &\psi'(0+)-\psi'(0-)=\alpha\psi(0)\},
 \end{split}
\end{equation}
and $H_\alpha\psi=-\psi''$ off zero.  This normalization is the owner;
changing the jump sign exchanges attractive and repulsive chambers.

Point interactions and their complete solvability are classical.  The
monograph of Albeverio, Gesztesy, H\o egh-Krohn, and Holden contains a direct
chapter on the one-center one-dimensional delta interaction
\cite{AGHH1988}.  Our contribution is the self-contained convention and
executable closure below, not a priority claim.

\section{One scalar equation gives the resolvent}
For $\kappa>0$, the free kernel at energy $-\kappa^2$ is
\[
 G_\kappa(x,y)=\frac{e^{-\kappa|x-y|}}{2\kappa}.
\]

\begin{theorem}[Rank-one resolvent]\label{thm:resolvent}
If $2\kappa+\alpha\ne0$, then
\begin{equation}\label{eq:resolvent}
 (H_\alpha+\kappa^2)^{-1}(x,y)
 =\frac1{2\kappa}\left\{e^{-\kappa|x-y|}
 -\frac{\alpha}{2\kappa+\alpha}
 e^{-\kappa(|x|+|y|)}\right\}.
\end{equation}
\end{theorem}
\begin{proof}
Away from $x=y,0$, both exponentials solve the homogeneous equation.  The
free term is continuous and its derivative has jump $-1$ at $x=y$.
At zero, the kernel value in~\eqref{eq:resolvent} is
$e^{-\kappa|y|}/(2\kappa+\alpha)$.  The image term has derivative jump
$\alpha e^{-\kappa|y|}/(2\kappa+\alpha)$, exactly $\alpha$ times that value.
It therefore satisfies~\eqref{eq:domain} and the resolvent source equation.
Uniqueness follows from invertibility away from the displayed pole.
\end{proof}

The formula is a Krein rank-one correction.  The exceptional equation
$2\kappa+\alpha=0$ has a positive solution precisely when $\alpha<0$; it is
spectral data, not a missing regular grid point.

\ifnum\CRevisionRound>0
\section{Spectrum and two-channel scattering}
The odd subspace automatically vanishes at zero, so $H_\alpha$ is free there.
On the even subspace the interface becomes the half-line Robin condition
$\psi'(0+)=(\alpha/2)\psi(0)$.  The odd sine transform and the even Robin
generalized eigenfunctions
\[
 \varphi_{\alpha,k}(x)=\sqrt{\frac2\pi}\,
 \frac{2k\cos(kx)+\alpha\sin(kx)}{\sqrt{4k^2+\alpha^2}},
 \qquad x>0,\qquad k>0,
\]
give the two continuum transforms; when $\alpha<0$, the even transform is
completed by the single normalized bound vector in~\eqref{eq:bound}.

\begin{theorem}[Complete sign atlas]\label{thm:spectrum}
The essential spectrum is $[0,\infty)$ and is purely absolutely continuous;
the singular-continuous spectrum is empty.  If $\alpha<0$, there is exactly
one simple eigenvalue and normalized eigenfunction,
\begin{equation}\label{eq:bound}
 E_b=-\frac{\alpha^2}{4},\qquad
 \phi_b(x)=\sqrt{-\frac\alpha2}\,e^{\alpha|x|/2}.
\end{equation}
For $\alpha\ge0$ there is no eigenvalue.

For momentum $k>0$, a unit-amplitude wave incident from the left has
\begin{equation}\label{eq:scattering}
 r_\alpha(k)=\frac{\alpha}{2ik-\alpha},\qquad
 t_\alpha(k)=\frac{2ik}{2ik-\alpha}.
\end{equation}
The scattering matrix has odd eigenvalue $1$ and even eigenvalue
$(2ik+\alpha)/(2ik-\alpha)$, both of modulus one.
\end{theorem}
\begin{proof}
The resolvent pole gives~\eqref{eq:bound}; direct integration gives
$\|\phi_b\|_2^2=(-\alpha/2)(-2/\alpha)=1$.  Stone's formula applied to the
explicit odd Dirichlet and even Robin resolvents yields the Lebesgue-density
transforms displayed above.  Their denominator $4k^2+\alpha^2$ has no
positive-real zero, the zero-energy solution is not in $L^2$, and the only
off-continuum singularity is the displayed attractive pole.  Hence there is
no embedded point or singular-continuous spectrum.  For scattering, write
$e^{ikx}+re^{-ikx}$ on $x<0$ and $te^{ikx}$ on $x>0$.  Continuity gives
$t=1+r$ and the jump gives $2ikr=\alpha t$, proving~\eqref{eq:scattering}.
Finally
\[
 |r|^2=\frac{\alpha^2}{4k^2+\alpha^2},\qquad
 |t|^2=\frac{4k^2}{4k^2+\alpha^2},
\]
so flux is one; diagonalizing the symmetric left--right matrix gives the two
channel phases.
\end{proof}

At $\alpha=0$ the operator is exactly free.  For nonzero coupling, the
zero-energy limit is perfectly reflecting and the high-energy limit is
transparent.  The attractive pole occurs at $k=i|\alpha|/2$; a repulsive
coupling has no upper-half-plane pole.
\fi

\ifnum\CRevisionRound>1
\section{Heat dynamics and the relative trace}
Let $K_0(t;u)=(4\pi t)^{-1/2}e^{-u^2/(4t)}$.  Laplace inversion of the two
terms in~\eqref{eq:resolvent} uses, for
$\Re\sqrt{s}>\max\{0,-\alpha/2\}$, the one-line identity
\[
 \mathcal L_{t\to s}\!\left[
 e^{\alpha a/2+\alpha^2t/4}
 \erfc\!\left(\frac{a}{2\sqrt t}+\frac{\alpha\sqrt t}{2}\right)\right]
 =\frac{2e^{-a\sqrt{s}}}{\sqrt{s}\,(2\sqrt{s}+\alpha)}.
\]
It yields the following all-time expression.

\begin{theorem}[Heat kernel and trace defect]\label{thm:heat}
For $t>0$, $x,y\in\R$, and $a=|x|+|y|$,
\begin{equation}\label{eq:heat}
 K_\alpha(t;x,y)=K_0(t;x-y)-\frac\alpha4
 e^{\alpha a/2+\alpha^2t/4}
 \erfc\!\left(\frac{a}{2\sqrt t}+\frac{\alpha\sqrt t}{2}\right).
\end{equation}
The diagonal difference from the free heat kernel is integrable and
\begin{equation}\label{eq:trace}
 \int_\R\{K_\alpha(t;x,x)-K_0(t;0)\}\,\D x
 =\frac12\left[e^{\alpha^2t/4}
 \erfc\!\left(\frac{\alpha\sqrt t}{2}\right)-1\right].
\end{equation}
\end{theorem}
\begin{proof}
The Laplace transform in $t$ of~\eqref{eq:heat} is
\eqref{eq:resolvent}; uniqueness of the semigroup transform proves the kernel.
For the trace, symmetry reduces the image term to
\[
 -\frac\alpha2e^{\alpha^2t/4}
 \int_0^\infty e^{\alpha x}
 \erfc\!\left(\frac{x}{\sqrt t}+\frac{\alpha\sqrt t}{2}\right)\D x.
\]
One integration by parts evaluates the integral as
\[
 \frac{e^{-\alpha^2t/4}-
 \erfc(\alpha\sqrt t/2)}{\alpha},
\]
with its continuous value at $\alpha=0$, giving~\eqref{eq:trace}.
\end{proof}

The trace defect tends to zero as $t\downarrow0$.  For $\alpha<0$, its
large-time leading term is $e^{\alpha^2t/4}$, the heat weight of the unique
negative eigenvalue.  This is not growth of the unitary group
$e^{-itH_\alpha}$, which remains norm preserving.  For $\alpha>0$, the
relative trace tends to $-1/2$.

\section{Executable closure and Route-A boundary}
The producer contains 32 regular resolvent cells, three deliberately
separated poles, 28 scattering cells, three bound-state cells, and eight
high-precision heat cells.  A strict duplicate-rejecting checker reports
1,726 assertions and reconstructs the heat values by inverse Laplace
transformation and an integrated diagonal resolvent; a symbolic engine
reports 46 identities; two fresh outputs are byte identical; and 30/30
semantic, structural, raw-JSON, and stale-hash attacks fail as required.
These finite receipts audit signs and factors; they do not prove
Theorems~\ref{thm:resolvent}--\ref{thm:heat}.

The natural self-adjoint Hamiltonian earns only
\texttt{A4\_NATURAL\_QUANTIZATION}.  No rational-prime carrier, primitive
repetition law, logarithmic arithmetic clock, target divisor, target
functional equation, or Hilbert--P\'olya operator is constructed.  Under the
literal scope \texttt{NO\_BAD\_EULER\_OR\_ROOT\_NUMBER}, the strict tuple is
\[
 \begin{gathered}
 \texttt{(A0\_FAIL,A1\_FAIL,A2\_FAIL,}\\[-2pt]
 \texttt{A3\_FAIL,A4\_NATURAL\_QUANTIZATION)},
 \end{gathered}
\]
the overall verdict is \texttt{ROUTE\_A\_REJECTED}, and Route B is disabled.
\fi

\begin{thebibliography}{1}
\bibitem{AGHH1988}
S. Albeverio, F. Gesztesy, R. H\o egh-Krohn, and H. Holden,
\emph{Solvable Models in Quantum Mechanics}, Theoretical and Mathematical
Physics, Springer, 1988,
\href{https://doi.org/10.1007/978-3-642-88201-2}
{doi:10.1007/978-3-642-88201-2}.
\end{thebibliography}
\end{document}
